\section{Why the Two Proofs Look Different}
\label{sec:imp-vs-lam}

We have now mechanized two Hoare-style soundness theorems against
the same propositional forcing semantics: the \texttt{while} rule
of Section~\ref{sec:imp} and the \texttt{wp\_fix} rule of
Section~\ref{sec:lam}.  Both are statements about recursively
defined programs.  Both go by some form of recursion in the proof.
But the proofs look structurally different, and that difference is
the main technical observation we want to articulate.

The IMP proof of \texttt{hoare\_while} proceeds by an explicit outer
induction on the step index $n$.  The base case discharges the wp at
depth $0$ (which is \texttt{True}).  The inductive step unfolds the
wp at depth $n+1$, case-analyses on the step taken by \texttt{CWhile b
c}, and in the recursive branch uses the induction hypothesis at
depth $n$ to discharge the obligation for \texttt{CWhile b c} itself.
The proof can also be done by an appeal to the internal Löb rule, as
the \texttt{hoare\_while\_iLob} alternative shows; but the two
proofs are structurally identical, and the choice between them is
cosmetic.

The $\lambda$-calculus proof of \texttt{wp\_fix} also goes by some
form of recursion --- it has to, because the goal is the soundness
of a recursive function --- but it does not contain any induction on
$n$ at all.  It factors the goal through Löb's form
$\triangleright \varphi \to \varphi$ and lets the recursion happen
inside \texttt{iLob}.  The interior of the factored proof is finite,
non-recursive step-indexed reasoning: a single case analysis on the
step index supplied by the \texttt{iImpl}, one unfolding of the
fix-application step via \texttt{wp\_pure\_step\_at}, and one
appeal to the Lipschitz hypothesis on the body.

\paragraph{Why does the difference arise?}
The structural reason is what we call a difference in the
\emph{recursion locus} of the soundness predicate.

For \texttt{while}, the wp predicate
$\mathtt{wp\_at}\, n\, (\mathtt{CWhile}\, b\, c)\, Q\, s$
unfolds at depth $n+1$ into an obligation that mentions
$\mathtt{wp\_at}\, n\, (\mathtt{CWhile}\, b\, c)\, Q\, s'$ at depth
$n$ for some state $s'$.  The proposition recursive at the lower
depth is \emph{the very same proposition}, decremented in the index.
At the meta level, this is just induction on $n$: each level of the
recursion consumes one step index and we run out of indices in
finitely many steps.  No specifically modal reasoning is needed
because the recursive obligation is parameterised solely by $n$.

For \texttt{fix}, the picture is different.  The proposition
$\mathtt{recursive\_spec}\, (\mathtt{EFix}\, body)\, S$ at depth
$n+1$ unfolds, after one fix-unfolding step, into an obligation
about the same $\mathtt{EFix}\, body$ at depth $n$ ---
\emph{but only because the body's behaviour at the recursive call site is
guaranteed by the same proposition at depth $n$.}  In other words,
the recursive proposition appears both inside and outside the
recursion: as the spec we are trying to prove for $\mathtt{EFix}\, body$, and as the hypothesis we need at the recursive call.  At the
meta level, an outer induction on $n$ can in principle reach the
recursive call --- one can always recurse on $n$ --- but the proof
obligation at the recursive call is no longer a statement purely
about smaller step indices; it is a statement about
$\mathtt{EFix}\, body$ itself at a smaller step index.  This is
exactly the shape that Löb's rule was designed to handle:
``assume the spec holds later, prove it holds now'', where ``later''
means ``at a strictly smaller condition''.  Internalizing the
recursion into Löb's rule is what frees the proof from carrying an
outer induction on $n$.

\paragraph{A diagnostic.}
We can state this as a working slogan:
\begin{quote}
\textbf{If the soundness predicate is itself recursive in the
program} --- if its definition at one depth refers to its definition
for the same program at a lower depth --- then internal Löb is the
natural soundness tool, and a meta-level induction on the step index
fights the structure of the predicate.  If the soundness predicate
is only recursive in the depth, an outer induction is interchangeable
with Löb and either is cheap.
\end{quote}
\noindent
The IMP \texttt{wp\_at}, $\mathtt{hoare}$, and \texttt{while} fall on
the easy side of the diagnostic.  The recursive-spec
$\mathtt{recursive\_spec}\, (\mathtt{EFix}\, body)\, S$ falls on the
hard side, and Löb earns its keep there.

\paragraph{What this says about step-indexed Hoare logic.}
In the verification literature, Löb's rule is sometimes presented as
a convenience: a tactic available in the Iris proof mode, used when
recursion is involved, but conceptually equivalent to ``the
step-indexed argument''.  Our diagnostic says this is right for
first-order recursion (the IMP case) but understates what is
happening in higher-order recursion (the $\lambda$-fix case).  When
the soundness predicate is genuinely self-referential in the
program, an outer induction on $n$ is still possible but is no
longer playing to the structure of the predicate.  The internal Löb
rule is what plays to that structure --- and it is, at bottom,
well-founded induction on the forcing order, internalized into the
logic so that it can act on the recursive proposition directly.
Section~\ref{sec:abstract} will make precise what
``internalized'' means by showing the abstract framework's Löb
rule is literally an appeal to \texttt{well\_founded\_ind}.

